%% GENERATED by tools/gen_latex.py -- edit content/ch15.py instead.
%% Build:  latexmk -pdf ch15.tex     (or: pdflatex ch15.tex, twice)
\documentclass[aspectratio=169,11pt,t]{beamer}
\usepackage{itbpro}

\coursetitle{Language Models and the Transformer}
\coursesubtitle{From predicting one character of Shakespeare to fine-tuning a 124-million-parameter pretrained encoder --- and the one mechanism that made the difference.}
\coursekicker{Chapter 15}
\coursesource{Chollet \& Watson, Deep Learning with Python 3e --- chapter 15}
\courseduration{4 hours (3 sessions)}
\coursepresenter{Prof. Bambang Riyanto Trilaksono $\cdot$ Rahman Indra Kesuma, S.Kom., M.Cs.}{Lead Instructor $\cdot$ Teaching Assistant}
\coursebrand{AI for Professional \textperiodcentered\ ITB \textperiodcentered\ Ch 15}

\title{Language Models and the Transformer}
\author{Prof. Bambang Riyanto Trilaksono $\cdot$ Rahman Indra Kesuma, S.Kom., M.Cs.}
\date{}

\begin{document}
\covertitle

\framekicker{Learning outcomes}
\begin{frame}[shrink=25]{Session Objectives}
\leadpar{By the end of this session, participants can:}
\begin{isteps}
\item Define a \textbf{language model} as p(token | past tokens), and explain why predicting one token at a time makes text generation tractable at all.
\item Build and \textbf{sample from} an autoregressive character-level model, including the inference loop that has no counterpart in training.
\item Build \textbf{sequence-to-sequence} translation with RNNs, and state the two limits that killed the approach.
\item Derive \textbf{attention} from the problem it solves --- score, softmax, weighted sum --- and read query/key/value fluently.
\item Assemble the \textbf{Transformer encoder and decoder} blocks, and explain every part: residuals, layer normalization, causal masking, feedforward.
\item Explain why the Transformer \textbf{needs positional embeddings}, and what happens numerically when you leave them out.
\item \textbf{Fine-tune a pretrained RoBERTa} on a classification task with KerasHub.
\item Explain what attention is doing \textbf{geometrically}, and why interpolation produces both generalization and hallucination.
\end{isteps}

\vskip 6pt
\vskip 4pt
\reslink{Book}{Chapter 15 --- full text}{https://deeplearningwithpython.io/chapters/chapter15\_language-models-and-the-transformer}
\reslink{Notebook}{01\_shakespeare\_language\_model.ipynb}{https://hendrikarisma.my.id/rs/ai-products-course/notebooks/ch15/01\_shakespeare\_language\_model.html}
\reslink{Notebook}{02\_seq2seq\_rnn\_translation.ipynb}{https://hendrikarisma.my.id/rs/ai-products-course/notebooks/ch15/02\_seq2seq\_rnn\_translation.html}
\reslink{Notebook}{03\_attention\_from\_scratch.ipynb}{https://hendrikarisma.my.id/rs/ai-products-course/notebooks/ch15/03\_attention\_from\_scratch.html}
\reslink{Notebook}{04\_transformer\_translation.ipynb}{https://hendrikarisma.my.id/rs/ai-products-course/notebooks/ch15/04\_transformer\_translation.html}
\reslink{Notebook}{05\_finetuning\_roberta.ipynb}{https://hendrikarisma.my.id/rs/ai-products-course/notebooks/ch15/05\_finetuning\_roberta.html}
\reslink{Notebook}{All chapter 15 notebooks}{https://hendrikarisma.my.id/rs/ai-products-course/notebooks/index.html\#ch15}
\reslink{Repo}{Author's official notebook}{https://github.com/fchollet/deep-learning-with-python-notebooks/blob/master/chapter15\_language-models-and-the-transformer.ipynb}
\vskip 2pt
\end{frame}

\coursesection{01}{The language model}{One idea, small enough to fit on a line: p(token | past tokens).}

\framekicker{Section 15.1}
\begin{frame}[shrink=25]{Classification was the easy case}
\leadpar{The previous chapter classified movie reviews. Text went \textbf{in}; a single floating-point number came \textbf{out}.}
{\fontsize{9.4}{12.6}\selectfont\color{fg2}That output shape is unusually forgiving. Binary classification needs one number. \textit{N}-way classification needs \textit{N}. Both are small, fixed, and known in advance.\par}\vskip 4pt
{\fontsize{9.4}{12.6}\selectfont\color{fg2}Question answering and translation are not like that. They need a model that \hilite{generates text as output}. Just as we needed tokenizers and embeddings to handle text on the way \textbf{in}, we need new machinery to produce it on the way \textbf{out}.\par}\vskip 4pt
\begin{iband}[signal]
We do not start from scratch. An integer sequence is still the natural numeric representation for text --- we \textbf{detokenize} by running the mapping in reverse.
\end{iband}
\end{frame}
\note{Set up the asymmetry deliberately. Students who did chapter 14 will assume text generation is a small extension of text classification. It is not --- the output space is the whole problem.}

\framekicker{Section 15.1}
\begin{frame}[shrink=25]{The obvious approach does not survive arithmetic}
{\fontsize{9.4}{12.6}\selectfont\color{fg2}The simplest option is to train a direct classifier over the space of \textbf{all possible output sequences}. Let us count how many classes that would be.\par}\vskip 4pt
\begin{minipage}[t]{0.3253\textwidth}
\istat{20,000}{vocabulary size}
\end{minipage}%
\hfill
\begin{minipage}[t]{0.3253\textwidth}
\istat{160 quadrillion}{possible 4-word sequences}
\end{minipage}%
\hfill
\begin{minipage}[t]{0.3253\textwidth}
\istat{> atoms in the universe}{possible 20-word sequences}
\end{minipage}%
{\fontsize{9.4}{12.6}\selectfont\color{fg2}Twenty thousand to the fourth power is 1.6 $\times$ 10$^{1}$$^{7}$. No model design rescues this --- the output layer alone would overwhelm any compute budget that will ever exist.\par}\vskip 4pt
\begin{center}
\includegraphics[width=0.92\linewidth,height=0.60\textheight,keepaspectratio]{figs/ch15-intractable.tex.pdf}
\end{center}
\vskip 3pt{\color{ink3}\fontsize{7.4}{9.4}\selectfont Two ways to produce a sequence, and only one of them is finite.\par}
\end{frame}
\note{Have them do the multiplication out loud. The point is not the number, it is that the number ends the discussion --- this is a case where back-of-the-envelope arithmetic eliminates an entire design.}

\framekicker{Section 15.1 $\cdot$ definition}
\begin{frame}[shrink=25]{A language model predicts one token at a time}
\leadpar{A \textbf{language model} learns a straightforward but deep probability distribution: \hilite{p(token | past tokens)}}
{\fontsize{9.4}{12.6}\selectfont\color{fg2}Given every token observed up to a point, it outputs a probability distribution over all tokens that could come next. With a 20,000-word vocabulary, the model needs to produce \textbf{20,000 numbers} --- not 160 quadrillion.\par}\vskip 4pt
{\fontsize{9.4}{12.6}\selectfont\color{fg2}And yet, by predicting the next token repeatedly, feeding each prediction back in as input, we have built a model that can generate text of \hilite{any length at all}.\par}\vskip 4pt
\begin{center}
\includegraphics[width=0.92\linewidth,height=0.60\textheight,keepaspectratio]{figs/ch15-lm-loop.tex.pdf}
\end{center}
\vskip 3pt{\color{ink3}\fontsize{7.4}{9.4}\selectfont The autoregressive loop: the output at one step is the input at the next.\par}
\end{frame}
\note{This slide is the whole chapter in miniature, and arguably the whole second half of the book. Everything from here --- the Transformer, chapter 16's LLMs, chapter 17's diffusion analogues --- is engineering on top of this one definition.}

\framekicker{Section 15.1.1 $\cdot$ listing 15.1}
\begin{frame}[shrink=25]{A concrete case: Shakespeare, one character at a time}
{\fontsize{9.4}{12.6}\selectfont\color{fg2}To make the idea concrete we train the smallest useful language model: one that predicts the \textbf{next character}. The training corpus is a collection of Shakespeare's plays and sonnets.\par}\vskip 4pt
\icode{python}{listing 15.1}{listings/ch15-01.py}
{\fontsize{9.4}{12.6}\selectfont\color{fg2}A character-level model keeps the whole example small enough to train in about two minutes, while demonstrating every mechanism a billion-parameter model uses.\par}\vskip 4pt
\end{frame}

\framekicker{Section 15.1.1}
\begin{frame}[shrink=25]{What the corpus looks like}
{\fontsize{9.4}{12.6}\selectfont\color{fg2}Before designing anything, look at the data --- the habit from chapter 6 applies here exactly as it did to images.\par}\vskip 4pt
\ioutput{listings/ch15-02.txt}
{\fontsize{9.4}{12.6}\selectfont\color{fg2}Note what the model will have to learn: speaker names in capitals, a colon, a newline, then verse. \textbf{None of that is given to it.} It is all structure to be discovered from next-character prediction alone.\par}\vskip 4pt
\end{frame}
\note{Point at the structure explicitly. When the generated sample appears later with correct speaker-name formatting, students should remember that nobody encoded that rule.}

\framekicker{Section 15.1.1 $\cdot$ listing 15.2}
\begin{frame}[shrink=25]{Features and labels differ by exactly one character}
{\fontsize{9.4}{12.6}\selectfont\color{fg2}We chunk the text into equal-length pieces, as we did with weather measurements in chapter 13. Because the tokenizer is character-level, we can chunk the \textbf{raw string} directly.\par}\vskip 4pt
\icode{python}{listing 15.2}{listings/ch15-03.py}
\begin{iband}[signal]
\inlinecode{shakespeare[:-1]} against \inlinecode{shakespeare[1:]} --- the labels are the features \textbf{shifted one step into the future}. That single line is the entire supervision signal.
\end{iband}
\end{frame}

\framekicker{Section 15.1.1}
\begin{frame}[shrink=25]{The offset, seen in one sample}
{\fontsize{9.4}{12.6}\selectfont\color{fg2}Print an (x, y) pair and the shift is visible directly.\par}\vskip 4pt
\ioutput{listings/ch15-04.txt}
{\fontsize{9.4}{12.6}\selectfont\color{fg2}At every position in the sequence, the label is the \textbf{next} character. One 100-character input yields \hilite{100 supervised predictions}, not one --- which is why this trains quickly despite the tiny corpus.\par}\vskip 4pt
\end{frame}

\framekicker{Section 15.1.1 $\cdot$ listing 15.3}
\begin{frame}[shrink=25]{A character-level vocabulary needs 67 entries}
{\fontsize{9.4}{12.6}\selectfont\color{fg2}\inlinecode{TextVectorization} from chapter 14 does the job, with \inlinecode{split="character"} instead of the default whitespace splitting, and no standardization --- case and punctuation are passed through unaltered.\par}\vskip 4pt
\icode{python}{listing 15.3}{listings/ch15-05.py}
\ioutput{listings/ch15-06.txt}
{\fontsize{9.4}{12.6}\selectfont\color{fg2}Sixty-seven symbols cover the complete works. Compare the 20,000-word vocabulary of chapter 14 --- this is the \textbf{vocabulary/sequence-length trade-off} in its extreme form.\par}\vskip 4pt
\end{frame}

\framekicker{Section 15.1.1 $\cdot$ listing 15.4}
\begin{frame}[shrink=25]{The model: embed, recur, project}
{\fontsize{9.4}{12.6}\selectfont\color{fg2}An RNN is the natural fit here: its recurrent state carries information about past characters forward to the prediction of the current one. A single GRU keeps the model small.\par}\vskip 4pt
\icode{python}{listing 15.4}{listings/ch15-07.py}
{\fontsize{9.4}{12.6}\selectfont\color{fg2}\inlinecode{return\_sequences=True} is essential: we want a prediction at \textbf{every position}, not just at the end. The final \inlinecode{Dense} outputs a distribution over all 67 possible next characters.\par}\vskip 4pt
\end{frame}

\framekicker{Section 15.1.1}
\begin{frame}[shrink=25]{Four million parameters, and where they sit}
\ioutput{listings/ch15-08.txt}
{\fontsize{9.4}{12.6}\selectfont\color{fg2}Almost \textbf{98\% of the parameters are in the GRU}. In the Transformer models later in this chapter the balance shifts dramatically --- worth remembering the shape of this table.\par}\vskip 4pt
\end{frame}

\framekicker{Section 15.1.1 $\cdot$ listing 15.5}
\begin{frame}[shrink=25]{It is still classification --- 6,400 of them per batch}
{\fontsize{9.4}{12.6}\selectfont\color{fg2}Nothing about the compile step is new. The loss is ordinary categorical crossentropy; the optimizer is the Adam default from chapter 3.\par}\vskip 4pt
\icode{python}{listing 15.5}{listings/ch15-09.py}
{\fontsize{9.4}{12.6}\selectfont\color{fg2}The model still solves a classification problem --- it just makes one classification per \textbf{token}. A batch of 64 samples with 100 characters each produces \hilite{6,400 individual labels}.\par}\vskip 4pt
{\fontsize{9.4}{12.6}\selectfont\color{fg2}Keras averages loss and accuracy first across each sequence, then across the batch. After 20 epochs the model predicts the next character correctly about \textbf{70\%} of the time.\par}\vskip 4pt
\end{frame}
\note{70\% next-character accuracy sounds unimpressive until you note that English text at character level has an entropy such that this is quite good for a two-minute training run on a 1 MB corpus.}

\framekicker{Section 15.1.2}
\begin{frame}[shrink=25]{Generation needs surgery on the trained model}
\leadpar{A model called in a feedback loop, where its output at one step becomes its input at the next, is called \textbf{autoregressive}.}
\begin{minipage}[t]{0.494\textwidth}
\begin{icard}
\cardh{During training}
Fixed sequence length of 100 tokens. The GRU's state is handled \textbf{implicitly} inside the layer call.
\end{icard}
\end{minipage}%
\hfill
\begin{minipage}[t]{0.494\textwidth}
\begin{icardaccent}
\cardh{During generation}
One token at a time, and the GRU state must be \textbf{explicitly} returned so it can be passed back in on the next call.
\end{icardaccent}
\end{minipage}%
{\fontsize{9.4}{12.6}\selectfont\color{fg2}The computational structure is identical, so we build a second model with modified inputs and outputs and \hilite{copy the weights across}.\par}\vskip 4pt
\end{frame}

\framekicker{Section 15.1.2 $\cdot$ listing 15.6}
\begin{frame}[shrink=25]{The inference model, with state made explicit}
{\fontsize{9.4}{12.6}\selectfont\color{fg2}Two changes to the architecture: the input is a single token rather than a hundred, and the GRU's state is promoted from an internal detail to a \textbf{named input and a named output}.\par}\vskip 4pt
\icode{python}{listing 15.6}{listings/ch15-10.py}
{\fontsize{9.4}{12.6}\selectfont\color{fg2}\inlinecode{shape=(1,)} --- one character in. \inlinecode{return\_state=True} plus an \inlinecode{initial\_state} input --- the recurrent state now travels \textbf{through the function signature} rather than inside the layer.\par}\vskip 4pt
\end{frame}

\framekicker{Section 15.1.2 $\cdot$ listing 15.7}
\begin{frame}[shrink=25]{Priming: replaying a prompt to reach the right state}
{\fontsize{9.4}{12.6}\selectfont\color{fg2}Before generating we must \textbf{prime} the GRU with the prompt --- feed it in one character at a time so the state matches what the model would have seen had the prompt occurred during training.\par}\vskip 4pt
\icode{python}{listing 15.7}{listings/ch15-11.py}
{\fontsize{9.4}{12.6}\selectfont\color{fg2}When the last character of the prompt goes in, the returned state summarises the \hilite{entire prompt}, and the returned prediction is already the distribution for the first generated character.\par}\vskip 4pt
\end{frame}

\framekicker{Section 15.1.2 $\cdot$ listing 15.8}
\begin{frame}[shrink=25]{The sampling loop}
{\fontsize{9.4}{12.6}\selectfont\color{fg2}With a primed state and a first prediction in hand, generation is a loop over three operations --- and the whole loop lives outside the model.\par}\vskip 4pt
\icode{python}{listing 15.8}{listings/ch15-12.py}
{\fontsize{9.4}{12.6}\selectfont\color{fg2}Three lines of substance: take the most likely character, append it, feed it back. The state persists across iterations and carries everything the model knows about what it has already written.\par}\vskip 4pt
\end{frame}
\note{argmax here is the crudest possible sampling strategy. Chapter 16 replaces it with temperature sampling, top-k, and nucleus sampling --- flag that this is a placeholder.}

\framekicker{Section 15.1.2}
\begin{frame}[shrink=25]{What two minutes of training produces}
\ioutput{listings/ch15-13.txt}
{\fontsize{9.4}{12.6}\selectfont\color{fg2}Not the next great tragedy. But look at what a \textbf{next-character} objective produced: correctly spelled words, speaker names in capitals followed by a colon and a newline, blank lines between speeches, and roughly verse-length lines.\par}\vskip 4pt
\begin{iband}[amber]
We trained on the narrow problem of guessing one character. We are using it for the far broader problem of \hilite{open-ended text generation}. That gap is the point.
\end{iband}
\end{frame}

\framekicker{Section 15.1.2 $\cdot$ a diagnostic worth running}
\begin{frame}[shrink=25]{Replace GRU with Bidirectional(GRU) and watch it break}
{\fontsize{9.4}{12.6}\selectfont\color{fg2}This whole setup works \textbf{only because} a recurrent network passes information forward. Try the substitution and observe the failure mode.\par}\vskip 4pt
\begin{itable}{>{\raggedright\arraybackslash}p{0.3948\textwidth}>{\raggedright\arraybackslash}p{0.2632\textwidth}>{\raggedright\arraybackslash}p{0.282\textwidth}}
\thead{What you change} & \thead{Training accuracy} & \thead{Generation} \\ \midrule
\inlinecode{GRU(hidden\_dim, return\_sequences=True)} & $\approx$ 70\% after 20 epochs & Shakespeare-like text \\
\inlinecode{Bidirectional(GRU(...))} & \textbf{above 99\% immediately} & \textbf{stops working entirely} \\
\end{itable}
{\fontsize{9.4}{12.6}\selectfont\color{fg2}With a bidirectional layer, information from the \textit{next} token reaches the prediction of the current one. The model reads the label off its own input --- the problem becomes trivial, and nothing generalisable is learned.\par}\vskip 4pt
\begin{iband}[rose]
A training metric that suddenly looks \textit{too good} is the most reliable signal of leakage there is. Chapter 5's lesson, in a new costume.
\end{iband}
\end{frame}
\note{Make them actually run this in the notebook. Seeing 99\% accuracy paired with garbage output is worth more than any amount of explanation about causal masking later in the chapter.}

\framekicker{Section 15.1.2 $\cdot$ what is new here}
\begin{frame}[shrink=25]{There is now logic that exists only at inference time}
\leadpar{Every model so far in this book answered \inlinecode{model.predict()}. This one does not.}
\begin{ibullets}
\item There is an entire \textbf{loop}, and a non-trivial amount of logic, that exists \textbf{only at inference time}.
\item State looping inside the GRU cell happens during both training and inference --- that part is symmetric.
\item But at \textbf{no point during training} do we feed the model's own predicted labels back to it as input. That happens only when generating.
\end{ibullets}
{\fontsize{9.4}{12.6}\selectfont\color{fg2}This asymmetry --- training on ground-truth prefixes, generating from self-produced prefixes --- has a name in the literature (\textit{exposure bias}) and is one reason generated text degrades over long outputs.\par}\vskip 4pt
\end{frame}

\coursesection{02}{Sequence-to-sequence learning}{Two sequences, one model: translation, and the encoder/decoder template.}

\framekicker{Section 15.2}
\begin{frame}[shrink=25]{The encoder/decoder template}
{\fontsize{9.4}{12.6}\selectfont\color{fg2}\textbf{Sequence-to-sequence} (seq2seq) modelling takes a source text as fixed input and generates a target sequence. Translation is the classic case; question answering is another.\par}\vskip 4pt
\begin{minipage}[t]{0.494\textwidth}
\begin{icardaccent}
\cardh{Encoder}
Turns the \textbf{whole} source sequence into an intermediate representation. It may look forward and backward freely.
\end{icardaccent}
\end{minipage}%
\hfill
\begin{minipage}[t]{0.494\textwidth}
\begin{icardaccent}
\cardh{Decoder}
The language-model setup from section 15.1, with one addition: it also sees the encoder's representation of the source.
\end{icardaccent}
\end{minipage}%
\begin{center}
\includegraphics[width=0.92\linewidth,height=0.60\textheight,keepaspectratio]{figs/ch15-seq2seq.tex.pdf}
\end{center}
\vskip 3pt{\color{ink3}\fontsize{7.4}{9.4}\selectfont Figure 15.1 --- during training the decoder sees the true target prefix; during inference it sees only what it has generated so far.\par}
\end{frame}

\framekicker{Section 15.2.1}
\begin{frame}[shrink=25]{English to Spanish: 118,000 sentence pairs}
{\fontsize{9.4}{12.6}\selectfont\color{fg2}The dataset is a plain text file, one example per line: an English sentence, a tab, then the Spanish. Downloading and parsing it takes a dozen lines.\par}\vskip 4pt
\icode{python}{}{listings/ch15-14.py}
{\fontsize{9.4}{12.6}\selectfont\color{fg2}The one non-obvious line is the wrapping of every Spanish sentence in \inlinecode{[start]} and \inlinecode{[end]}.\par}\vskip 4pt
\end{frame}

\framekicker{Section 15.2.1}
\begin{frame}[shrink=25]{The seed and the stop signal live in the data}
{\fontsize{9.4}{12.6}\selectfont\color{fg2}Look at what one parsed pair contains.\par}\vskip 4pt
\ioutput{listings/ch15-15.txt}
{\fontsize{9.4}{12.6}\selectfont\color{fg2}\inlinecode{[start]} is what we feed the decoder to begin generating; \inlinecode{[end]} is what tells the generation loop to stop. They are \hilite{inserted in the data, not built into the model} --- which is why the tokenizer must be told not to strip square brackets.\par}\vskip 4pt
{\fontsize{9.4}{12.6}\selectfont\color{fg2}Splitting is the usual three-way shuffle: 70\% train, 15\% validation, 15\% test, taken over the pairs rather than over the sentences, so no English sentence appears in two splits.\par}\vskip 4pt
\end{frame}

\framekicker{Section 15.2.1 $\cdot$ listing 15.9}
\begin{frame}[shrink=25]{Two tokenizers, because punctuation is language-specific}
{\fontsize{9.4}{12.6}\selectfont\color{fg2}Two customisations are needed, and both are the kind of detail that silently ruins a translation model.\par}\vskip 4pt
\begin{ibullets}
\item The characters \inlinecode{[} and \inlinecode{]} are stripped by default. We must \textbf{keep} them, so the token \inlinecode{"[start]"} stays distinct from the word \inlinecode{"start"}.
\item Spanish uses \inlinecode{¿}, which is not in \inlinecode{string.punctuation}. The Spanish tokenizer needs it added to the strip set explicitly.
\end{ibullets}
\icode{python}{listing 15.9}{listings/ch15-16.py}
\begin{iband}[signal]
For a non-toy translator you would treat punctuation as \textbf{separate tokens} rather than stripping it --- otherwise the model can never produce correctly punctuated output.
\end{iband}
\end{frame}

\framekicker{Section 15.2.1 $\cdot$ listing 15.9}
\begin{frame}[shrink=25]{The two vectorizers, and why their lengths differ}
{\fontsize{9.4}{12.6}\selectfont\color{fg2}With the standardization function in place, the two \inlinecode{TextVectorization} layers are near-identical --- except for one number that is easy to miss.\par}\vskip 4pt
\icode{python}{listing 15.9 (cont.)}{listings/ch15-17.py}
{\fontsize{9.4}{12.6}\selectfont\color{fg2}The Spanish sequence is \hilite{one token longer}: 20 tokens of input and 20 of label, cut from 21. And \inlinecode{adapt()} runs on \inlinecode{train\_pairs} only --- never on validation or test text.\par}\vskip 4pt
\end{frame}

\framekicker{Section 15.2.1 $\cdot$ listing 15.10}
\begin{frame}[shrink=25]{Shifting the target, and masking the padding}
{\fontsize{9.4}{12.6}\selectfont\color{fg2}The pipeline must return \inlinecode{(inputs, target, sample\_weights)}, where \inlinecode{inputs} is a dictionary with two keys --- the tokenized English and the tokenized Spanish.\par}\vskip 4pt
\icode{python}{listing 15.10}{listings/ch15-18.py}
{\fontsize{9.4}{12.6}\selectfont\color{fg2}\inlinecode{spa[:, :-1]} is the decoder's input; \inlinecode{spa[:, 1:]} is its label. The same tensor, offset by one --- exactly the Shakespeare setup, with the encoder input added alongside.\par}\vskip 4pt
\begin{iband}[signal]
\inlinecode{sample\_weights = labels != 0} tells Keras to \textbf{ignore padded positions} when computing loss and metrics. Without it, a model that learns only to predict padding would score well.
\end{iband}
\end{frame}

\framekicker{Section 15.2.1}
\begin{frame}[shrink=25]{What comes out of the pipeline}
\ioutput{listings/ch15-19.txt}
{\fontsize{9.4}{12.6}\selectfont\color{fg2}Four aligned tensors of the same shape. Two go in as a dictionary of named inputs, one is the label, one is the per-position weight. Every seq2seq pipeline in this chapter has this shape.\par}\vskip 4pt
\end{frame}

\framekicker{Section 15.2.2}
\begin{frame}[shrink=25]{Why the naive single-RNN approach cannot work}
{\fontsize{9.4}{12.6}\selectfont\color{fg2}The easiest thing to try is one RNN reading source tokens and emitting target tokens at each step:\par}\vskip 4pt
\icode{python}{}{listings/ch15-20.py}
{\fontsize{9.4}{12.6}\selectfont\color{fg2}Because of the step-by-step nature of RNNs, this model sees only source tokens 0\ldots{}N when predicting target token N.\par}\vskip 4pt
\begin{iquote}
\quotetext{\textit{"I will bring the bag to you"} becomes \textit{"Te traeré la bolsa."} The \textbf{first} Spanish word corresponds to the \textbf{last} English word. There is no way to emit it without having read to the end of the source.}\quotecite{Section 15.2.2}
\end{iquote}
\end{frame}
\note{Any bilingual student will have an example of their own here --- Indonesian to English word-order shifts work just as well. Ask for one.}

\framekicker{Section 15.2.2 $\cdot$ figure 15.2}
\begin{frame}[shrink=25]{So read the whole source first}
{\fontsize{9.4}{12.6}\selectfont\color{fg2}A human translator reads the entire source sentence before beginning. The proper seq2seq setup does the same: an encoder RNN compresses the whole source into a \textbf{single vector}, which becomes the decoder's initial state.\par}\vskip 4pt
\begin{center}
\includegraphics[width=0.92\linewidth,height=0.60\textheight,keepaspectratio]{figs/ch15-rnn-bottleneck.tex.pdf}
\end{center}
\vskip 3pt{\color{ink3}\fontsize{7.4}{9.4}\selectfont Figure 15.2 --- everything the decoder knows about English arrives through one 1,024-dimensional vector.\par}
{\fontsize{9.4}{12.6}\selectfont\color{fg2}Instead of the zero initial state used in the Shakespeare generator, the decoder starts from a state that \hilite{encodes the source sentence}.\par}\vskip 4pt
\end{frame}

\framekicker{Section 15.2.2 $\cdot$ listings 15.11-15.12}
\begin{frame}[shrink=25]{Encoder bidirectional, decoder emphatically not}
{\fontsize{9.4}{12.6}\selectfont\color{fg2}Both halves are a few lines of Keras, but one design decision differs between them and it is the decisive one.\par}\vskip 4pt
\icode{python}{listing 15.11 - encoder}{listings/ch15-21.py}
\icode{python}{listing 15.12 - decoder}{listings/ch15-22.py}
\begin{iband}[signal]
\textbf{Bidirectional in the encoder is a good idea} --- we never predict source tokens, so there is nothing to cheat at. \textbf{Bidirectional in the decoder would break training} for exactly the reason the Shakespeare experiment showed.
\end{iband}
\end{frame}

\framekicker{Section 15.2.2}
\begin{frame}[shrink=25]{Thirty-five million parameters, and where they went}
\ioutput{listings/ch15-23.txt}
{\fontsize{9.4}{12.6}\selectfont\color{fg2}The output \inlinecode{Dense} alone is 15 million parameters --- 1,024 hidden units times a 15,000-word vocabulary. \textbf{The vocabulary projection dominates}, which is a constant of language modelling at every scale.\par}\vskip 4pt
\end{frame}

\framekicker{Section 15.2.2}
\begin{frame}[shrink=25]{65\% --- and why that number should not be trusted}
{\fontsize{9.4}{12.6}\selectfont\color{fg2}Training is unremarkable --- \inlinecode{weighted\_metrics} rather than \inlinecode{metrics} so that the sample weights we built actually reach the accuracy calculation.\par}\vskip 4pt
\icode{python}{}{listings/ch15-24.py}
{\fontsize{9.4}{12.6}\selectfont\color{fg2}We reach \textbf{65\% next-token accuracy}. But next-token accuracy is a poor metric for translation, and it is worth being precise about why.\par}\vskip 4pt
\begin{ibullets}
\item It assumes the correct target tokens 0\ldots{}N are \textbf{already known} when predicting token N+1. At inference, they are not --- you generate them.
\item It punishes a perfectly good translation that happens to use a different word order from the reference.
\item The standard alternative is \textbf{BLEU}, which compares against a set of high-quality reference translations and tolerates misalignment.
\end{ibullets}
\end{frame}
\note{This connects straight back to chapter 6's rule about choosing a metric that measures success rather than what is convenient to compute.}

\framekicker{Section 15.2.2 $\cdot$ listing 15.13}
\begin{frame}[shrink=25]{What the RNN translator actually produces}
\ioutput{listings/ch15-25.txt}
{\fontsize{9.4}{12.6}\selectfont\color{fg2}Decent for a toy, and still making basic mistakes --- the first example has the adverb wedged into the middle of the verb phrase.\par}\vskip 4pt
\begin{iband}[amber]
This inference loop is also \textbf{inefficient}: it reprocesses the entire source and the entire generated target every time it samples one new word. A real system caches the state that has not changed.
\end{iband}
\end{frame}

\framekicker{Section 15.2.2 $\cdot$ the limit}
\begin{frame}[shrink=25]{Two limits that no amount of tuning removes}
{\fontsize{9.4}{12.6}\selectfont\color{fg2}You can stack deeper recurrent layers, swap GRU for LSTM, widen the state. None of it addresses the following.\par}\vskip 4pt
\begin{minipage}[t]{0.494\textwidth}
\begin{icardbad}
\cardh{The bottleneck}
The \textbf{entire} source representation must fit in the encoder's state vector, which caps the size and complexity of sentence you can translate at all.
\end{icardbad}
\end{minipage}%
\hfill
\begin{minipage}[t]{0.494\textwidth}
\begin{icardbad}
\cardh{Forgetting}
RNNs progressively forget the past. By the 100th token, little information about the start of the sequence remains.
\end{icardbad}
\end{minipage}%
{\fontsize{9.4}{12.6}\selectfont\color{fg2}Recurrent networks dominated seq2seq in the mid-2010s --- \textbf{Google Translate circa 2017 was a stack of seven large LSTM layers}, essentially the model we just built. These two limits are what drove the search for something else.\par}\vskip 4pt
\end{frame}

\coursesection{03}{The Transformer architecture}{Attention is derived from the problem it solves, then made into everything.}

\framekicker{Section 15.3 $\cdot$ 2017}
\begin{frame}[shrink=25]{The finding is in the title}
\begin{iquote}
\quotetext{\textbf{Attention Is All You Need}}\quotecite{Vaswani et al., NeurIPS 2017 --- arxiv.org/abs/1706.03762}
\end{iquote}
{\fontsize{9.4}{12.6}\selectfont\color{fg2}The authors were working on translation systems like the one we just built. Attention itself was not new --- it had been in NLP systems for a couple of years. The \textbf{surprise} was that it was useful enough to be the \textit{only} mechanism passing information across a sequence.\par}\vskip 4pt
{\fontsize{9.4}{12.6}\selectfont\color{fg2}No recurrent layers at all. This finding unleashed a revolution in natural language processing, and then beyond it --- vision, audio, protein structure, and the models in chapters 16 and 17.\par}\vskip 4pt
\end{frame}

\framekicker{Section 15.3 $\cdot$ the intuition}
\begin{frame}[shrink=25]{A thought experiment about this book}
\leadpar{Suppose you had to build a weather-prediction model using only this textbook.}
\begin{ibullets}
\item You might read the whole book cover to cover --- over 100,000 words, far longer than any sequence we have handled.
\item But when you actually wrote the code, you would pay \textbf{special attention} to the timeseries chapter, and within it to a few specific listings.
\item You would not be worried about the details of image convolutions.
\end{ibullets}
{\fontsize{9.4}{12.6}\selectfont\color{fg2}Humans are \textbf{selective and contextual} about where they pull information from. An RNN has no mechanism to refer back directly --- all information must pass through the cell state, in a loop, through every intervening position.\par}\vskip 4pt
\begin{iband}[signal]
It is like finishing the book, closing it, and writing the weather model \hilite{entirely from memory}.
\end{iband}
\end{frame}

\framekicker{Section 15.3.1}
\begin{frame}[shrink=25]{Score every source vector against the current target}
{\fontsize{9.4}{12.6}\selectfont\color{fg2}Return to the translation RNN. When predicting a single token we have one vector for the target position, and a \textbf{sequence} of vectors for the source words.\par}\vskip 4pt
{\fontsize{9.4}{12.6}\selectfont\color{fg2}The goal: give the model a way to score every source vector by its relevance to the word being predicted. Assume for a moment a function \inlinecode{score(target\_vector, source\_vector)}.\par}\vskip 4pt
\begin{center}
\includegraphics[width=0.92\linewidth,height=0.60\textheight,keepaspectratio]{figs/ch15-attention-idea.tex.pdf}
\end{center}
\vskip 3pt{\color{ink3}\fontsize{7.4}{9.4}\selectfont Figure 15.4 --- attention assigns a relevance score to each source vector, for each target vector.\par}
\end{frame}

\framekicker{Section 15.3.1}
\begin{frame}[shrink=25]{The scoring function is a dot product, and the sum is a softmax}
{\fontsize{9.4}{12.6}\selectfont\color{fg2}Two design choices, both made for the simplest available reason. \textbf{Score} by dot product --- vectors close in an embedding space score high. \textbf{Normalise} with a softmax, so the scores for a given target sum to 1 and the weighted sum has predictable magnitude.\par}\vskip 4pt
\icode{python}{}{listings/ch15-26.py}
{\fontsize{9.4}{12.6}\selectfont\color{fg2}The \inlinecode{einsum} subscripts spell out the shapes: \inlinecode{b}atch, \inlinecode{t}arget length, \inlinecode{s}ource length, \inlinecode{d}imension. The first contraction produces a \textbf{(batch, target, source) score matrix}; the second uses it to take a weighted sum.\par}\vskip 4pt
\end{frame}
\note{If einsum is unfamiliar, write the two operations out as loops on the board once. Every attention implementation in every framework is these two contractions with a softmax between them.}

\framekicker{Section 15.3.1 $\cdot$ figure 15.5}
\begin{frame}[shrink=25]{The score matrix is two-dimensional}
{\fontsize{9.4}{12.6}\selectfont\color{fg2}When both target and source are sequences, the attention scores form a \textbf{matrix}: one row per target word, one column per source word.\par}\vskip 4pt
\begin{itable}{>{\raggedright\arraybackslash}p{0.2068\textwidth}>{\raggedright\arraybackslash}p{0.1034\textwidth}>{\raggedright\arraybackslash}p{0.1034\textwidth}>{\raggedright\arraybackslash}p{0.1034\textwidth}>{\raggedright\arraybackslash}p{0.1034\textwidth}>{\raggedright\arraybackslash}p{0.1034\textwidth}>{\raggedright\arraybackslash}p{0.1034\textwidth}>{\raggedright\arraybackslash}p{0.1128\textwidth}}
\thead{target $\downarrow$ / source $\rightarrow$} & \thead{I} & \thead{will} & \thead{bring} & \thead{the} & \thead{bag} & \thead{to} & \thead{you} \\ \midrule
Te & $\cdot$ & $\cdot$ & $\cdot$ & $\cdot$ & $\cdot$ & $\cdot$ & \textbf{0.8} \\
traeré & 0.2 & \textbf{0.5} & \textbf{0.3} & $\cdot$ & $\cdot$ & $\cdot$ & $\cdot$ \\
la & $\cdot$ & $\cdot$ & $\cdot$ & \textbf{0.6} & 0.3 & $\cdot$ & $\cdot$ \\
bolsa & $\cdot$ & $\cdot$ & $\cdot$ & 0.2 & \textbf{0.7} & $\cdot$ & $\cdot$ \\
{}[end] & $\cdot$ & $\cdot$ & $\cdot$ & $\cdot$ & $\cdot$ & 0.3 & 0.4 \\
\end{itable}
{\fontsize{9.4}{12.6}\selectfont\color{fg2}Read a row as: \textit{when producing this Spanish word, how much did the model draw on each English word?} The \inlinecode{Te} row peaking at \inlinecode{you} is exactly the long-range dependency the RNN could not express.\par}\vskip 4pt
\end{frame}

\framekicker{Section 15.3.1}
\begin{frame}[shrink=25]{Give the model parameters: query, key, value}
{\fontsize{9.4}{12.6}\selectfont\color{fg2}The hypothesis space gets much richer if we project both source and target with \inlinecode{Dense} layers first, letting the model find a shared space in which relevant vectors are close.\par}\vskip 4pt
\icode{python}{}{listings/ch15-27.py}
{\fontsize{9.4}{12.6}\selectfont\color{fg2}\inlinecode{sum(score(target, source) * source)} has become \inlinecode{sum(score(query, key) * value)}. The three-argument form is more general: in rare cases you want different vectors for \textbf{scoring} and for \textbf{summing}.\par}\vskip 4pt
\end{frame}

\framekicker{Section 15.3.1 $\cdot$ figure 15.6}
\begin{frame}[shrink=25]{The names come from search engines}
{\fontsize{9.4}{12.6}\selectfont\color{fg2}Imagine a tool that looks up photos in a database.\par}\vskip 4pt
\begin{minipage}[t]{0.3253\textwidth}
\begin{icardaccent}
\cardh{Query}
Your search term. In attention: the target vector for the position currently being decoded.
\end{icardaccent}
\end{minipage}%
\hfill
\begin{minipage}[t]{0.3253\textwidth}
\begin{icardaccent}
\cardh{Keys}
The photo tags used to match against the query. In attention: the projected source vectors that get scored.
\end{icardaccent}
\end{minipage}%
\hfill
\begin{minipage}[t]{0.3253\textwidth}
\begin{icardaccent}
\cardh{Values}
The photos themselves --- what you actually retrieve. In attention: the projected source vectors that get summed.
\end{icardaccent}
\end{minipage}%
\begin{center}
\includegraphics[width=0.92\linewidth,height=0.60\textheight,keepaspectratio]{figs/ch15-search.tex.pdf}
\end{center}
\vskip 3pt{\color{ink3}\fontsize{7.4}{9.4}\selectfont Figure 15.6 --- the query is compared to keys, and the match scores rank the values.\par}
\end{frame}

\framekicker{Section 15.3.1 $\cdot$ refinement 1}
\begin{frame}[shrink=25]{Scale the scores before the softmax}
{\fontsize{9.4}{12.6}\selectfont\color{fg2}When input vectors get long, dot products get large. Large logits push the softmax toward a one-hot output, where its \textbf{gradients vanish} --- training becomes unstable.\par}\vskip 4pt
{\fontsize{9.4}{12.6}\selectfont\color{fg2}The fix is one division:\par}\vskip 4pt
\icode{python}{}{listings/ch15-28.py}
{\fontsize{9.4}{12.6}\selectfont\color{fg2}Scaling by the square root of the vector length works well for any vector size --- the variance of a dot product of \textit{d}-dimensional vectors grows with \textit{d}, so dividing by $\sqrt{}$\textit{d} holds it roughly constant. This is why the mechanism is called \hilite{scaled dot-product attention}.\par}\vskip 4pt
\end{frame}

\framekicker{Section 15.3.1 $\cdot$ refinement 2}
\begin{frame}[shrink=25]{One softmax sum is too blunt}
{\fontsize{9.4}{12.6}\selectfont\color{fg2}The weighted sum is powerful --- a direct connection across distant parts of a sequence. It is also \textbf{blunt}: if the model attends to many tokens at once, the interesting features of individual source tokens get washed out in the combined representation.\par}\vskip 4pt
{\fontsize{9.4}{12.6}\selectfont\color{fg2}The trick that works: run the whole attention operation \textbf{several times in parallel}, with different parameters. Each run is an \textit{attention head}.\par}\vskip 4pt
\begin{iband}[signal]
By projecting query and key differently, one head might learn to match the \textbf{subject} of the source sentence, while another attends to \textbf{punctuation}.
\end{iband}
\begin{center}
\includegraphics[width=0.92\linewidth,height=0.74\textheight,keepaspectratio]{figs/ch15-multihead.tex.pdf}
\end{center}
\vskip 3pt{\color{ink3}\fontsize{7.4}{9.4}\selectfont Figure 15.7 --- each head attends to different parts of the source, in separate partitions of the eventual output vector.\par}
\end{frame}

\framekicker{Section 15.3.1}
\begin{frame}[shrink=25]{Self-attention, computed on one sentence}
\begin{center}
\includegraphics[width=0.97\linewidth,height=0.56\textheight,keepaspectratio]{figs/ch15-qkv-run.tex.pdf}
\end{center}
\vskip 3pt{\color{ink3}\fontsize{7.4}{9.4}\selectfont One query --- the word ``it'' --- against every key in the sentence. Press play to walk the five stages.\par}
\end{frame}

\framekicker{Section 15.3.1}
\begin{frame}[shrink=25]{Reading the five stages}
\begin{isteps}
\item \textbf{Three matrices, one embedding.} \inlinecode{Wq}, \inlinecode{Wk} and \inlinecode{Wv} turn each token's embedding into a query, a key and a value. They are learned, and they are the \textit{only} parameters in an attention head.
\item \textbf{The query is dotted against every key.} That gives one raw score per token --- how well ``it'' matches each word, including itself. Dividing by \inlinecode{sqrt(d)} keeps the scores in a range softmax can use.
\item \textbf{Softmax turns scores into proportions.} They are positive and they sum to 100\%. Nothing is discarded: every token gets a share, and a low share is still a share.
\item \textbf{The output is the value vectors, mixed in those proportions.} \hilite{The new vector for ``it'' is literally built out of the other words in the sentence.} That is the whole mechanism.
\end{isteps}
\begin{iband}[signal]
This is one \textbf{head}. Multi-head attention runs several of these in parallel with different \inlinecode{Wq}, \inlinecode{Wk}, \inlinecode{Wv}, so each can specialise --- one on syntax, another on which noun a pronoun refers to --- and concatenates the results.
\end{iband}
\end{frame}

\framekicker{Section 15.3.1 $\cdot$ figure 15.7}
\begin{frame}[shrink=25]{Multi-head attention, written out}
{\fontsize{9.4}{12.6}\selectfont\color{fg2}Every projection becomes a \textbf{list} of projections, one per head, and the per-head outputs are concatenated at the end.\par}\vskip 4pt
\icode{python}{}{listings/ch15-29.py}
{\fontsize{9.4}{12.6}\selectfont\color{fg2}Each head produces a \inlinecode{head\_dim}-wide output; concatenating \inlinecode{num\_heads} of them and projecting once gives the layer's output. In practice the loop is one batched matmul.\par}\vskip 4pt
\end{frame}

\framekicker{Section 15.3.1}
\begin{frame}[shrink=25]{In Keras, this is one layer}
{\fontsize{9.4}{12.6}\selectfont\color{fg2}Everything on the previous three slides is a reusable layer, and Keras ships it.\par}\vskip 4pt
\icode{python}{}{listings/ch15-30.py}
{\fontsize{9.4}{12.6}\selectfont\color{fg2}We derived it by hand first for a reason: \inlinecode{MultiHeadAttention} is the single most consequential layer in modern deep learning, and reading its arguments --- \inlinecode{query}, \inlinecode{key}, \inlinecode{value}, \inlinecode{attention\_mask}, \inlinecode{use\_causal\_mask} --- is a skill the rest of this course assumes.\par}\vskip 4pt
\end{frame}

\framekicker{Section 15.3.2 $\cdot$ self-attention}
\begin{frame}[shrink=25]{Now point the mechanism at its own input}
{\fontsize{9.4}{12.6}\selectfont\color{fg2}So far attention has passed information \textbf{between} two sequences. But nothing stops us from letting a sequence attend to itself:\par}\vskip 4pt
\icode{python}{}{listings/ch15-31.py}
{\fontsize{9.4}{12.6}\selectfont\color{fg2}This is \textbf{self-attention}. Each token attends to every token in its own sequence, including itself, producing a representation of the word \hilite{in context}.\par}\vskip 4pt
\begin{iquote}
\quotetext{\textit{"The train left the station on time."} What kind of station? A radio station? The International Space Station? Self-attention lets the model give a high score to the pair (\textbf{station}, \textbf{train}), summing the representation of \textit{train} into that of \textit{station}.}\quotecite{Section 15.3.2}
\end{iquote}
\end{frame}
\note{This is the slide where the architecture stops being a translation trick and becomes a general representation-learning mechanism. Every encoder model in chapters 15-17 is built on this line of code.}

\framekicker{Section 15.3.2 $\cdot$ the missing ingredient}
\begin{frame}[shrink=25]{Attention alone collapses to a matrix multiplication}
{\fontsize{9.4}{12.6}\selectfont\color{fg2}Self-attention sounds like what an RNN does. Can we simply swap the recurrent layers for \inlinecode{MultiHeadAttention}? \textbf{Almost --- but not quite.}\par}\vskip 4pt
{\fontsize{9.4}{12.6}\selectfont\color{fg2}\inlinecode{MultiHeadAttention} combines \textit{linear} projections of source elements. That is all it does. It is a very expressive \textbf{pooling} operation, and pooling is not enough.\par}\vskip 4pt
\begin{iband}[rose]
Consider a sequence of length one. The score matrix is a single 1, and the layer reduces to a linear projection. You could stack \textbf{100} such layers and still simplify the whole computation to \hilite{one matrix multiplication}.
\end{iband}
{\fontsize{9.4}{12.6}\selectfont\color{fg2}Every recurrent cell eventually passes each token's vector through a dense projection with an activation. The Transformer adds this back in the simplest possible way: a feedforward network of \textbf{two Dense layers with a nonlinearity between them}.\par}\vskip 4pt
\end{frame}

\framekicker{Section 15.3.2}
\begin{frame}[shrink=25]{The division of labour inside a Transformer block}
\begin{minipage}[t]{0.494\textwidth}
\begin{icardaccent}
\cardh{Attention}
Passes information \textbf{across} the sequence. Mixes positions. Contains no nonlinearity.
\end{icardaccent}
\end{minipage}%
\hfill
\begin{minipage}[t]{0.494\textwidth}
\begin{icardaccent}
\cardh{Feedforward}
Updates the representation of \textbf{each item independently}. Mixes features, not positions. Contains the nonlinearity.
\end{icardaccent}
\end{minipage}%
{\fontsize{9.4}{12.6}\selectfont\color{fg2}Two more ingredients come straight from the ConvNet chapters: \textbf{residual connections} and \textbf{normalization}, both of which chapter 9 established as essential to training anything deep.\par}\vskip 4pt
\begin{center}
\includegraphics[width=0.92\linewidth,height=0.60\textheight,keepaspectratio]{figs/ch15-encoder.tex.pdf}
\end{center}
\vskip 3pt{\color{ink3}\fontsize{7.4}{9.4}\selectfont The encoder block: attend, add, norm; feed forward, add, norm.\par}
\end{frame}

\framekicker{Section 15.3.2 $\cdot$ listing 15.14}
\begin{frame}[shrink=25]{The encoder block: what it holds}
{\fontsize{9.4}{12.6}\selectfont\color{fg2}Five sub-layers: one attention, two dense, two normalization.\par}\vskip 4pt
\icode{python}{listing 15.14 - constructor}{listings/ch15-32.py}
{\fontsize{9.4}{12.6}\selectfont\color{fg2}\inlinecode{key\_dim = hidden\_dim // num\_heads} keeps the total width constant as heads are added --- eight heads of 32 dimensions rather than one head of 256. \textbf{Splitting, not growing.}\par}\vskip 4pt
\end{frame}

\framekicker{Section 15.3.2 $\cdot$ listing 15.14}
\begin{frame}[shrink=25]{The encoder block: what it does}
{\fontsize{9.4}{12.6}\selectfont\color{fg2}Two stages, each of the form \textit{transform, add the residual, normalise} --- the shape of every Transformer block ever written.\par}\vskip 4pt
\icode{python}{listing 15.14 - call}{listings/ch15-33.py}
{\fontsize{9.4}{12.6}\selectfont\color{fg2}Input and output have the \textbf{same shape}, so blocks stack --- each one building a progressively more expressive representation of the source sentence. RoBERTa, later in this chapter, stacks twelve.\par}\vskip 4pt
\end{frame}

\framekicker{Section 15.3.2 $\cdot$ a deliberate choice}
\begin{frame}[shrink=25]{LayerNormalization, not BatchNormalization}
{\fontsize{9.4}{12.6}\selectfont\color{fg2}The normalization here is \textbf{not} the \inlinecode{BatchNormalization} used in image models. Write both out and the reason is immediate.\par}\vskip 4pt
\icode{python}{layer normalization}{listings/ch15-34.py}
\icode{python}{batch normalization}{listings/ch15-35.py}
{\fontsize{9.4}{12.6}\selectfont\color{fg2}\inlinecode{BatchNormalization} pools over \textbf{axis 0}, creating interactions between samples in a batch. \inlinecode{LayerNormalization} pools only over the last axis, normalising each sequence \hilite{independently} --- which is what variable-length sequence data requires.\par}\vskip 4pt
\end{frame}
\note{Ask why interactions between samples are harmful here specifically. Answer: sequences in a batch have wildly different lengths and padding amounts, so batch statistics are contaminated by how you happened to group examples.}

\framekicker{Section 15.3.2 $\cdot$ masking}
\begin{frame}[shrink=25]{The padding mask, and its shape}
{\fontsize{9.4}{12.6}\selectfont\color{fg2}\inlinecode{attention\_mask} is a Boolean tensor broadcast to the shape of the attention scores --- \textbf{(batch\_size, target\_length, source\_length)}. Where it is false, the score is zeroed and that source token contributes nothing.\par}\vskip 4pt
{\fontsize{9.4}{12.6}\selectfont\color{fg2}We use it to stop any token from attending to \textbf{padding}, which carries no information.\par}\vskip 4pt
\icode{python}{}{listings/ch15-36.py}
{\fontsize{9.4}{12.6}\selectfont\color{fg2}The \inlinecode{None} in the middle is the whole trick: one mask, broadcast across every row of the score matrix, because \textit{which source tokens are padding} does not depend on which target position is asking.\par}\vskip 4pt
\end{frame}

\framekicker{Section 15.3.3}
\begin{frame}[shrink=25]{The decoder attends twice}
{\fontsize{9.4}{12.6}\selectfont\color{fg2}The decoder block is almost identical to the encoder, with one addition: it must also see the \textbf{encoder's output}. So it uses attention twice.\par}\vskip 4pt
\begin{isteps}
\item \textbf{Self-attention} over the target sequence, so each target position can use information from other target positions --- masked so it can only look backward.
\item \textbf{Cross-attention}, whose query is the target and whose key and value are the encoder output --- this is what brings information across from the source language.
\item \textbf{Feedforward}, same as the encoder.
\end{isteps}
\begin{center}
\includegraphics[width=0.92\linewidth,height=0.60\textheight,keepaspectratio]{figs/ch15-decoder.tex.pdf}
\end{center}
\vskip 3pt{\color{ink3}\fontsize{7.4}{9.4}\selectfont The decoder block: masked self-attention, then cross-attention, then feedforward --- each with add-and-norm.\par}
\end{frame}

\framekicker{Section 15.3.3 $\cdot$ listing 15.15}
\begin{frame}[shrink=25]{The decoder block: what it holds}
{\fontsize{9.4}{12.6}\selectfont\color{fg2}The constructor is the encoder's, plus a second \inlinecode{MultiHeadAttention} and its normalization layer.\par}\vskip 4pt
\icode{python}{listing 15.15 - constructor}{listings/ch15-37.py}
{\fontsize{9.4}{12.6}\selectfont\color{fg2}Nine attributes, three of them \inlinecode{LayerNormalization} --- one per residual junction. Counting normalization layers is a quick way to check a Transformer implementation by eye.\par}\vskip 4pt
\end{frame}

\framekicker{Section 15.3.3 $\cdot$ listing 15.15}
\begin{frame}[shrink=25]{The decoder block: what it does}
{\fontsize{9.4}{12.6}\selectfont\color{fg2}Three stages, each of the form \textit{transform, add the residual, normalise}. The two attention calls receive \textbf{different masks}.\par}\vskip 4pt
\icode{python}{listing 15.15 - call}{listings/ch15-38.py}
{\fontsize{9.4}{12.6}\selectfont\color{fg2}\inlinecode{use\_causal\_mask=True} on self-attention stops the decoder seeing its own future. The padding \inlinecode{attention\_mask} on cross-attention stops it attending to empty source positions. \textbf{Different problems, different masks.}\par}\vskip 4pt
\end{frame}

\framekicker{Section 15.3.3 $\cdot$ the causal mask}
\begin{frame}[shrink=25]{Attention is bidirectional by default --- which would be fatal}
{\fontsize{9.4}{12.6}\selectfont\color{fg2}In self-attention, \textbf{any} position can attend to \textbf{any} other. Without special care, the decoder would see the token it is trying to predict --- the same cheat as the \inlinecode{Bidirectional(GRU)} experiment in section 15.1.\par}\vskip 4pt
{\fontsize{9.4}{12.6}\selectfont\color{fg2}The fix is a lower-triangular attention mask:\par}\vskip 4pt
\ioutput{listings/ch15-39.txt}
{\fontsize{9.4}{12.6}\selectfont\color{fg2}Row \textit{i} is the mask for target position \textit{i}. Row 0: the first token may attend only to itself. Row 1: the second may attend to the first and itself. And so on --- the same forward-only information flow an RNN gets for free from its structure.\par}\vskip 4pt
\end{frame}

\framekicker{Section 15.3.3 $\cdot$ the causal mask}
\begin{frame}[shrink=25]{One argument, and a bug you will not be warned about}
{\fontsize{9.4}{12.6}\selectfont\color{fg2}Read the mask row by row and the rule is plain: position \textit{i} may look at everything up to and including itself, and nothing after.\par}\vskip 4pt
\begin{center}
\includegraphics[width=0.92\linewidth,height=0.60\textheight,keepaspectratio]{figs/ch15-causal.tex.pdf}
\end{center}
\vskip 3pt{\color{ink3}\fontsize{7.4}{9.4}\selectfont The same forward-only information flow an RNN gets free from its structure, imposed here by hand.\par}
\begin{iband}[rose]
In Keras this is one argument: \inlinecode{use\_causal\_mask=True}. \hilite{Forgetting it is one of the most common and most silent bugs in hand-written Transformer code}.
\end{iband}
\end{frame}

\framekicker{Section 15.3.4 $\cdot$ listing 15.16}
\begin{frame}[shrink=25]{Assembling the translator: the encoder half}
{\fontsize{9.4}{12.6}\selectfont\color{fg2}Same structure as the RNN model, GRU layers replaced. The feedforward block scales \textbf{up} to 2048 before the nonlinearity and back down to 256 after --- a large intermediate dimension that works well in practice.\par}\vskip 4pt
\icode{python}{listing 15.16 - encoder}{listings/ch15-40.py}
{\fontsize{9.4}{12.6}\selectfont\color{fg2}\inlinecode{source != 0} computes the padding mask inline --- no separate layer, just a Boolean tensor built from the input.\par}\vskip 4pt
\end{frame}

\framekicker{Section 15.3.4 $\cdot$ listing 15.16}
\begin{frame}[shrink=25]{Assembling the translator: the decoder half}
{\fontsize{9.4}{12.6}\selectfont\color{fg2}The decoder takes both the target embeddings and the encoder output, and the same source mask is handed to it again.\par}\vskip 4pt
\icode{python}{listing 15.16 - decoder}{listings/ch15-41.py}
{\fontsize{9.4}{12.6}\selectfont\color{fg2}One encoder block and one decoder block --- the 2017 paper used \textbf{six of each}. Everything else, down to the final softmax over the Spanish vocabulary, is unchanged from the GRU model.\par}\vskip 4pt
\end{frame}

\framekicker{Section 15.3.4}
\begin{frame}[shrink=25]{Fourteen million parameters --- less than half the RNN}
\ioutput{listings/ch15-42.txt}
{\fontsize{9.4}{12.6}\selectfont\color{fg2}34.9 M for the RNN against \textbf{14.4 M} here. The encoder and decoder blocks together are under 3 M parameters --- the embeddings and the vocabulary projection dominate, as always.\par}\vskip 4pt
\end{frame}

\framekicker{Section 15.3.4 $\cdot$ a deliberate failure}
\begin{frame}[shrink=25]{58\%. Worse than the RNN. Can you spot why?}
{\fontsize{9.4}{12.6}\selectfont\color{fg2}The compile and fit calls are identical to the RNN model's, so the two numbers are directly comparable.\par}\vskip 4pt
\icode{python}{}{listings/ch15-43.py}
\begin{minipage}[t]{0.494\textwidth}
\istat{65\%}{GRU seq2seq}
\end{minipage}%
\hfill
\begin{minipage}[t]{0.494\textwidth}
\istat{58\%}{Transformer, first attempt}
\end{minipage}%
{\fontsize{9.4}{12.6}\selectfont\color{fg2}Seven percentage points \textbf{worse}. Either the architecture is not what it was hyped up to be, or something is missing from the implementation.\par}\vskip 4pt
\begin{iband}[amber]
Stop here. Give the room two minutes with the code on the previous slide before turning over.
\end{iband}
\end{frame}
\note{Do not rush this. The reveal only lands if people have genuinely tried to find it. Hint if needed: 'this section is about sequence models --- is the model we just built a sequence model?'}

\framekicker{Section 15.3.4 $\cdot$ the reveal}
\begin{frame}[shrink=25]{What we built is not a sequence model at all}
\leadpar{It is composed of dense layers that process tokens independently, and an attention layer that looks at tokens \textbf{as a set}.}
{\fontsize{9.4}{12.6}\selectfont\color{fg2}Change the order of the tokens and you get identical pairwise attention scores and identical context-aware representations. Shuffle every word of every English source sentence and the model \hilite{would not notice}. You would get the same accuracy.\par}\vskip 4pt
\begin{center}
\includegraphics[width=0.92\linewidth,height=0.60\textheight,keepaspectratio]{figs/ch15-order-blind.tex.pdf}
\end{center}
\vskip 3pt{\color{ink3}\fontsize{7.4}{9.4}\selectfont Attention is a set-processing mechanism: it sees relationships between pairs, and is blind to where in the sequence they sit.\par}
{\fontsize{9.4}{12.6}\selectfont\color{fg2}For RNNs, order-awareness came free from the layer's computation. For the Transformer, we must \textbf{inject positional information into the embeddings themselves}.\par}\vskip 4pt
\end{frame}

\framekicker{Section 15.3.5}
\begin{frame}[shrink=25]{Positional embedding: add where to what}
{\fontsize{9.4}{12.6}\selectfont\color{fg2}The idea is simple. Each input embedding gets two components: the usual \textbf{word vector}, which represents the word independently of context, and a \textbf{position vector}, which represents where the word sits in this sentence.\par}\vskip 4pt
\begin{center}
\includegraphics[width=0.92\linewidth,height=0.60\textheight,keepaspectratio]{figs/ch15-posemb.tex.pdf}
\end{center}
\vskip 3pt{\color{ink3}\fontsize{7.4}{9.4}\selectfont Two embedding tables, added together. The model works out how to use the extra information.\par}
{\fontsize{9.4}{12.6}\selectfont\color{fg2}The most obvious scheme --- concatenate the raw integer position --- is a poor one: positions can be large integers, and neural networks dislike large input values and discrete input distributions.\par}\vskip 4pt
\end{frame}

\framekicker{Section 15.3.5 $\cdot$ positional encoding}
\begin{frame}[shrink=25]{The sinusoids, and the pattern they make}
\begin{center}
\includegraphics[width=0.97\linewidth,height=0.56\textheight,keepaspectratio]{figs/ch15-posenc-waves.tex.pdf}
\end{center}
\vskip 3pt{\color{ink3}\fontsize{7.4}{9.4}\selectfont Left: three of the actual waves. Right: every dimension for every position. Both computed from the formula, not drawn to look right.\par}
\end{frame}

\framekicker{Section 15.3.5 $\cdot$ two schemes}
\begin{frame}[shrink=25]{Sinusoidal, or learned}
\begin{minipage}[t]{0.494\textwidth}
\begin{icard}
\cardh{Sinusoidal (the 2017 paper)}
Add a vector of values in [-1, 1] that varies \textbf{cyclically} with position, built from cosine functions. Clever: it characterises any integer in a large range using only small values.
\end{icard}
\end{minipage}%
\hfill
\begin{minipage}[t]{0.494\textwidth}
\begin{icardaccent}
\cardh{Learned (what we use)}
Learn positional vectors the same way we learn word embeddings --- an \inlinecode{Embedding} table indexed by position. \textbf{Simpler and more effective} in practice.
\end{icardaccent}
\end{minipage}%
{\fontsize{9.4}{12.6}\selectfont\color{fg2}The learned scheme has one cost: it fixes a \textbf{maximum sequence length} at build time, since the table has one row per position. Sinusoidal encodings extrapolate; learned ones do not.\par}\vskip 4pt
\end{frame}

\framekicker{Section 15.3.5 $\cdot$ listing 15.17}
\begin{frame}[shrink=25]{The PositionalEmbedding layer}
{\fontsize{9.4}{12.6}\selectfont\color{fg2}Two \inlinecode{Embedding} tables in one layer: one indexed by token, one indexed by position. The \inlinecode{call()} adds them.\par}\vskip 4pt
\icode{python}{listing 15.17}{listings/ch15-44.py}
{\fontsize{9.4}{12.6}\selectfont\color{fg2}\inlinecode{ops.cumsum(ops.ones\_like(inputs), axis=-1) - 1} produces \inlinecode{[0, 1, 2, \ldots{}]} for every sequence in the batch --- a backend-agnostic way to write \inlinecode{arange} that works under TensorFlow, PyTorch, and JAX alike.\par}\vskip 4pt
{\fontsize{9.4}{12.6}\selectfont\color{fg2}The layer is a \textbf{drop-in replacement} for \inlinecode{Embedding}.\par}\vskip 4pt
\end{frame}

\framekicker{Section 15.3.5 $\cdot$ listing 15.18}
\begin{frame}[shrink=25]{Two lines changed}
{\fontsize{9.4}{12.6}\selectfont\color{fg2}Here is the whole model again. Compare it against listing 15.16 and find the difference before reading the note underneath.\par}\vskip 4pt
\icode{python}{listing 15.18}{listings/ch15-45.py}
{\fontsize{9.4}{12.6}\selectfont\color{fg2}Everything else is untouched. \inlinecode{layers.Embedding} became \inlinecode{PositionalEmbedding}, twice.\par}\vskip 4pt
\end{frame}

\framekicker{Section 15.3.5}
\begin{frame}[shrink=25]{67\% --- and a third of the epoch time}
\begin{minipage}[t]{0.3253\textwidth}
\istat{65\%}{GRU seq2seq, 34.9 M params}
\end{minipage}%
\hfill
\begin{minipage}[t]{0.3253\textwidth}
\istat{58\%}{Transformer without positions}
\end{minipage}%
\hfill
\begin{minipage}[t]{0.3253\textwidth}
\istat{67\%}{Transformer with positions, 14.4 M params}
\end{minipage}%
{\fontsize{9.4}{12.6}\selectfont\color{fg2}A noticeable improvement on the GRU --- and all the more so given that this model has \textbf{half the parameters}.\par}\vskip 4pt
{\fontsize{9.4}{12.6}\selectfont\color{fg2}There is a second thing to notice about the training run: each epoch takes about \hilite{one third} the time. With attention there is no looped state passing, so on a GPU or TPU the entire attention computation happens \textbf{in one go}.\par}\vskip 4pt
\begin{iband}[signal]
That speed-up is not a minor convenience. It is the property that made scaling to billions of parameters economically possible --- the subject of chapter 16.
\end{iband}
\end{frame}
\note{This is the causal chain worth stating out loud: parallelisable training $\rightarrow$ affordable scaling $\rightarrow$ large pretrained models $\rightarrow$ everything from 2018 onward. Not accuracy on this toy task.}

\framekicker{Section 15.3.5 $\cdot$ listing 15.19}
\begin{frame}[shrink=25]{The Transformer's translations}
\ioutput{listings/ch15-46.txt}
{\fontsize{9.4}{12.6}\selectfont\color{fg2}Subjectively much better than the GRU translations. It is still a toy model --- but it is a \textbf{better toy model}.\par}\vskip 4pt
\end{frame}

\framekicker{Section 15.3.5 $\cdot$ an honest caveat}
\begin{frame}[shrink=25]{None of these choices is provably optimal}
{\fontsize{9.4}{12.6}\selectfont\color{fg2}After all this implementation detail one might reasonably protest that it seems arbitrary. So many small decisions taken on faith. How could we know this configuration is optimal?\par}\vskip 4pt
\leadpar{\textbf{The answer is simple --- it is not.}}
\begin{ibullets}
\item Many improvements to attention, normalization, and positional embeddings have been proposed since 2017.
\item Much current research replaces attention with something \textbf{less computationally complex}, as sequence lengths grow very long.
\item Something will eventually supplant the Transformer --- possibly before you finish this course.
\end{ibullets}
\begin{iband}[signal]
The field moves \textbf{empirically}. Attention grew out of an attempt to augment RNNs; after years of guessing and checking by a great many people, it gave rise to the Transformer. There is little reason to think that process is done.
\end{iband}
\end{frame}

\coursesection{04}{Classification with a pretrained Transformer}{124 million parameters, one epoch of fine-tuning, and a three-point jump.}

\framekicker{Section 15.4}
\begin{frame}[shrink=25]{Two properties that invited scaling}
\begin{minipage}[t]{0.494\textwidth}
\begin{icardgood}
\cardh{Faster to train}
No loops during training, which is always good on a GPU or TPU. The whole sequence is processed at once.
\end{icardgood}
\end{minipage}%
\hfill
\begin{minipage}[t]{0.494\textwidth}
\begin{icardaccent}
\cardh{Data hungry}
We saw this ourselves: the RNN plateaued on validation after about 5 epochs; the Transformer was \textbf{still improving after 30}.
\end{icardaccent}
\end{minipage}%
{\fontsize{9.4}{12.6}\selectfont\color{fg2}These observations prompted many to try scaling the Transformer up --- more data, more layers, more parameters --- with strikingly good results.\par}\vskip 4pt
{\fontsize{9.4}{12.6}\selectfont\color{fg2}The result was a distinctive shift in the field toward \textbf{large pretrained models} that can cost millions to train but perform noticeably better across a wide range of text problems.\par}\vskip 4pt
\end{frame}

\framekicker{Section 15.4.1 $\cdot$ BERT}
\begin{frame}[shrink=25]{Masked language modelling}
{\fontsize{9.4}{12.6}\selectfont\color{fg2}\textbf{BERT} --- Bidirectional Encoder Representations from Transformers --- appeared a year after \textit{Attention Is All You Need}. Its structure was \textbf{exactly the encoder} of the translation Transformer we just built.\par}\vskip 4pt
{\fontsize{9.4}{12.6}\selectfont\color{fg2}But BERT was 100 to 300 million parameters, against our 14 million. That needs a great deal of training data, so the authors used a variant of the language-model objective.\par}\vskip 4pt
\begin{isteps}
\item Take a sequence of text and replace about \textbf{15\% of tokens} with a special \inlinecode{[MASK]} token.
\item Train the model to predict the \textbf{original} value of each masked token.
\item Note what is \textit{not} needed: \textbf{any labels at all}. For any text sequence you can choose random tokens and mask them.
\end{isteps}
\begin{iband}[signal]
A causal language model learns p(token | \textbf{past} tokens). A masked language model learns p(token | \textbf{surrounding} tokens).
\end{iband}
\end{frame}

\framekicker{Section 15.4.1 $\cdot$ the two families}
\begin{frame}[shrink=25]{Encoder or decoder --- a choice that follows from the objective}
\begin{center}
\includegraphics[width=0.92\linewidth,height=0.60\textheight,keepaspectratio]{figs/ch15-lm-kinds.tex.pdf}
\end{center}
\vskip 3pt{\color{ink3}\fontsize{7.4}{9.4}\selectfont The objective determines the masking, the masking determines what the model is good for.\par}
\begin{itable}{>{\raggedright\arraybackslash}p{0.2444\textwidth}>{\raggedright\arraybackslash}p{0.3478\textwidth}>{\raggedright\arraybackslash}p{0.3478\textwidth}}
\thead{} & \thead{Causal LM} & \thead{Masked LM} \\ \midrule
Distribution & p(token \textbackslash\{\}| past) & p(token \textbackslash\{\}| surrounding) \\
Attention & Causal-masked, one direction & Bidirectional \\
Good at & \textbf{Generating} text in a loop & \textbf{Representing} text richly \\
Examples & GPT, Llama --- chapter 16 & BERT, RoBERTa --- this section \\
\end{itable}
{\fontsize{9.4}{12.6}\selectfont\color{fg2}Pretrained word embeddings were already common practice when BERT appeared --- chapter 14 did exactly that. Pretraining a whole Transformer gave something more powerful: an embedding for a word \hilite{in the context of the words around it}.\par}\vskip 4pt
\end{frame}

\framekicker{Section 15.4.2 $\cdot$ RoBERTa}
\begin{frame}[shrink=25]{Same architecture, ten times the data}
{\fontsize{9.4}{12.6}\selectfont\color{fg2}We will use \textbf{RoBERTa} --- Robustly Optimized BERT --- rather than BERT itself. It made minor architectural simplifications, but the notable change was the training data.\par}\vskip 4pt
\begin{minipage}[t]{0.3253\textwidth}
\istat{16 GB}{BERT --- mainly Wikipedia}
\end{minipage}%
\hfill
\begin{minipage}[t]{0.3253\textwidth}
\istat{160 GB}{RoBERTa --- text from across the web}
\end{minipage}%
\hfill
\begin{minipage}[t]{0.3253\textwidth}
\istat{\textasciitilde{}\$300k}{estimated training cost at the time}
\end{minipage}%
{\fontsize{9.4}{12.6}\selectfont\color{fg2}Because of the extra data, RoBERTa performs noticeably better \textbf{at an equivalent parameter count}. That relationship --- data quantity buying quality at fixed model size --- is the central empirical fact of the pretraining era.\par}\vskip 4pt
\end{frame}

\framekicker{Section 15.4.2}
\begin{frame}[shrink=25]{Three things you need to use a pretrained model}
\begin{minipage}[t]{0.3253\textwidth}
\begin{icardaccent}
\cardh{A matching tokenizer}
Text must be tokenized \textbf{exactly} as it was during pretraining. If words map to different indices, the learned representations are meaningless.
\end{icardaccent}
\end{minipage}%
\hfill
\begin{minipage}[t]{0.3253\textwidth}
\begin{icardaccent}
\cardh{A matching architecture}
The internal maths must be recreated exactly. Ours nearly matches \inlinecode{TransformerEncoder} already --- but \textit{nearly} is not enough.
\end{icardaccent}
\end{minipage}%
\hfill
\begin{minipage}[t]{0.3253\textwidth}
\begin{icardaccent}
\cardh{The pretrained weights}
Produced by training for about a day on \textbf{1,024 GPUs} over billions of words.
\end{icardaccent}
\end{minipage}%
{\fontsize{9.4}{12.6}\selectfont\color{fg2}Recreating the first two by hand would be possible but time-consuming and error-prone. As in chapter 8, we use \textbf{KerasHub} instead.\par}\vskip 4pt
\end{frame}

\framekicker{Section 15.4.2 $\cdot$ listing 15.20}
\begin{frame}[shrink=25]{Two lines to load a pretrained Transformer}
{\fontsize{9.4}{12.6}\selectfont\color{fg2}KerasHub exposes pretrained checkpoints through a single constructor, \inlinecode{from\_preset()}, named by a string identifier.\par}\vskip 4pt
\icode{python}{listing 15.20}{listings/ch15-47.py}
{\fontsize{9.4}{12.6}\selectfont\color{fg2}\inlinecode{from\_preset()} loads weights, configuration, and tokenizer assets together --- which is precisely what keeps the three requirements on the previous slide consistent with one another.\par}\vskip 4pt
\ioutput{listings/ch15-48.txt}
{\fontsize{9.4}{12.6}\selectfont\color{fg2}RoBERTa's tokenizer is close to the subword tokenizer built in chapter 14, with tweaks to handle Unicode from any language.\par}\vskip 4pt
\end{frame}

\framekicker{Section 15.4.2 $\cdot$ why subwords are mandatory here}
\begin{frame}[shrink=25]{The tokenization decision, at 160 GB}
\begin{itable}{>{\raggedright\arraybackslash}p{0.2444\textwidth}>{\raggedright\arraybackslash}p{0.6956\textwidth}}
\thead{Choice} & \thead{What goes wrong at this scale} \\ \midrule
\textbf{Character level} & Input sequences become far too long, making the model much more expensive to train --- attention cost grows with the \textit{square} of length. \\
\textbf{Word level} & Covering the distinct words across millions of web documents blows up the vocabulary, making the front \inlinecode{Embedding} layer unworkably large. \\
\textbf{Subword (BPE)} & Handles \textbf{any} word with a 50,000-term vocabulary. This is why every large model uses it. \\
\end{itable}
{\fontsize{9.4}{12.6}\selectfont\color{fg2}The trade-off from chapter 14 has not changed --- but the pressure on both ends is far greater, which is what makes the middle option not merely convenient but \hilite{required}\par}\vskip 4pt
\end{frame}

\framekicker{Section 15.4.2 $\cdot$ what is a backbone}
\begin{frame}[shrink=25]{A model without a head}
{\fontsize{9.4}{12.6}\selectfont\color{fg2}Chapter 8 used \textit{backbone} for a vision network mapping images to a latent space --- a model without a prediction head. In KerasHub the term covers any pretrained model \textbf{not yet specialised for a task}.\par}\vskip 4pt
{\fontsize{9.4}{12.6}\selectfont\color{fg2}What we loaded maps an input sequence to an output sequence of shape \textbf{(batch\_size, sequence\_length, 768)}. It is not set up for any particular loss function.\par}\vskip 4pt
\begin{center}
\includegraphics[width=0.92\linewidth,height=0.60\textheight,keepaspectratio]{figs/ch15-backbone.tex.pdf}
\end{center}
\vskip 3pt{\color{ink3}\fontsize{7.4}{9.4}\selectfont One backbone, many heads: classifying sentences, extracting spans, tagging parts of speech.\par}
\begin{iband}[signal]
Think of it as attaching different heads to a screwdriver --- a Phillips head for one task, a flat head for another.
\end{iband}
\end{frame}

\framekicker{Section 15.4.2}
\begin{frame}[shrink=25]{124 million parameters, twelve blocks deep}
\ioutput{listings/ch15-49.txt}
{\fontsize{9.4}{12.6}\selectfont\color{fg2}Twelve identical encoder blocks stacked --- exactly the stacking property we noted when the \inlinecode{TransformerEncoder} preserved its input shape. This is the \textbf{smallest} RoBERTa checkpoint, and the largest model used anywhere in this book so far.\par}\vskip 4pt
\end{frame}

\framekicker{Section 15.4.3 $\cdot$ listing 15.21}
\begin{frame}[shrink=25]{Packing tokens the way pretraining did}
{\fontsize{9.4}{12.6}\selectfont\color{fg2}RoBERTa expects sequences packed in a specific form: an \inlinecode{<s>} token, the content, an \inlinecode{</s>} token, then \inlinecode{<pad>} tokens.\par}\vskip 4pt
\ioutput{listings/ch15-50.txt}
\icode{python}{listing 15.21}{listings/ch15-51.py}
{\fontsize{9.4}{12.6}\selectfont\color{fg2}Matching the pretraining token ordering as closely as possible makes the model train \textbf{faster and more accurately}. The IMDb loading code from chapter 14 is reused unchanged.\par}\vskip 4pt
\end{frame}

\framekicker{Section 15.4.3}
\begin{frame}[shrink=25]{One preprocessed batch}
\ioutput{listings/ch15-52.txt}
{\fontsize{9.4}{12.6}\selectfont\color{fg2}Batch size is now \textbf{16}, not the 32 or 64 used earlier --- a 124-million-parameter model with 512-token sequences is a different memory proposition. Rows ending in \inlinecode{1} are padded; the \inlinecode{padding\_mask} marks exactly where.\par}\vskip 4pt
\end{frame}

\framekicker{Section 15.4.3 $\cdot$ sidebar}
\begin{frame}[shrink=25]{Where does pretraining data come from?}
\begin{itable}{>{\raggedright\arraybackslash}p{0.2256\textwidth}>{\raggedright\arraybackslash}p{0.7144\textwidth}}
\thead{Model} & \thead{Pretraining data} \\ \midrule
The first Transformer (2017) & A well-known English-German translation set, \textbf{4 million sentence pairs}. \\
BERT (2018) & A dump of English Wikipedia plus a dataset of \textbf{7,000 self-published books}. \\
GPT-2 (2019) & Scraped by following \textbf{outgoing links from Reddit}. \\
Llama (latest) & "\textbf{15 trillion tokens} of data from publicly available sources." \\
\end{itable}
{\fontsize{9.4}{12.6}\selectfont\color{fg2}The short answer is: the internet. The other answer is that this has increasingly become a \textbf{secret} --- companies often do not release the data or the precise mixture of sources.\par}\vskip 4pt
\begin{iband}[amber]
When possible, pay close attention to where a model's data came from. It shapes both the \textbf{biases} and the \textbf{performance} of everything built on top of it.
\end{iband}
\end{frame}
\note{This sidebar is where governance, procurement, and legal questions enter the technical curriculum. Chapter 18 returns to it; chapter 16 shows how much the exact data mixture matters.}

\framekicker{Section 15.4.4 $\cdot$ listing 15.22}
\begin{frame}[shrink=25]{The classification head, and why it uses token zero}
{\fontsize{9.4}{12.6}\selectfont\color{fg2}The backbone outputs a whole sequence, shape (batch, length, 768). To predict one label we must condense that to a single vector per sample.\par}\vskip 4pt
\icode{python}{listing 15.22}{listings/ch15-53.py}
{\fontsize{9.4}{12.6}\selectfont\color{fg2}\inlinecode{x[:, 0, :]} takes the \textbf{first token's} representation. Mean or max pooling would also work, but this works slightly better --- and the reason is the nature of attention.\par}\vskip 4pt
\begin{iband}[signal]
The first position in the final encoder layer can attend to \textbf{every} other position and pull information from them. So rather than pooling with something coarse like an average, attention pools \hilite{contextually}.
\end{iband}
\end{frame}

\framekicker{Section 15.4.4 $\cdot$ listing 15.23}
\begin{frame}[shrink=25]{One epoch, and a very small learning rate}
{\fontsize{9.4}{12.6}\selectfont\color{fg2}One detail in this compile call carries almost all of the risk, and it is the optimizer's learning rate.\par}\vskip 4pt
\icode{python}{listing 15.23}{listings/ch15-54.py}
{\fontsize{9.4}{12.6}\selectfont\color{fg2}\inlinecode{5e-5} is roughly \textbf{twenty times smaller} than Adam's default. This is the fine-tuning discipline from chapter 8, unchanged: large updates would destroy the representations that cost \$300,000 to learn.\par}\vskip 4pt
\ioutput{listings/ch15-55.txt}
{\fontsize{9.4}{12.6}\selectfont\color{fg2}\textbf{93.7\% test accuracy after a single epoch}, against the 90\% ceiling reached in chapter 14.\par}\vskip 4pt
\end{frame}

\framekicker{Section 15.4.4 $\cdot$ the scoreboard}
\begin{frame}[shrink=25]{What the last two chapters bought, in one table}
\begin{itable}{>{\raggedright\arraybackslash}p{0.376\textwidth}>{\raggedright\arraybackslash}p{0.2444\textwidth}>{\raggedright\arraybackslash}p{0.3196\textwidth}}
\thead{Model} & \thead{IMDb test accuracy} & \thead{Cost} \\ \midrule
Bigram bag-of-words (chapter 14) & $\approx$ 89\% & seconds, on a CPU \\
Sequence model with pretrained embeddings & $\approx$ 90\% & minutes, one GPU \\
\textbf{RoBERTa base}, fine-tuned 1 epoch & \textbf{93.7\%} & 124 M params, one GPU-hour \\
RoBERTa large, fine-tuned & \textbf{> 95\%} & 300 M params \\
\end{itable}
\begin{center}
\includegraphics[width=0.92\linewidth,height=0.60\textheight,keepaspectratio]{figs/ch15-finetune.tex.pdf}
\end{center}
\vskip 3pt{\color{ink3}\fontsize{7.4}{9.4}\selectfont Where the accuracy actually came from: months of unlabelled pretraining, then one epoch on 20,000 labelled reviews.\par}
{\fontsize{9.4}{12.6}\selectfont\color{fg2}This is a far more expensive model to run than the bigram classifier, and the clear benefit has to be weighed against that. \textbf{Chapter 18 makes that trade-off explicit}; for now, note that the three points cost roughly four orders of magnitude in compute.\par}\vskip 4pt
\end{frame}

\coursesection{05}{What makes the Transformer effective?}{A geometric answer, by way of a model from 2013.}

\framekicker{Section 15.5 $\cdot$ 2013}
\begin{frame}[shrink=25]{Word2Vec, and the magic vectors}
{\fontsize{9.4}{12.6}\selectfont\color{fg2}At Google in 2013, Tomas Mikolov and colleagues built a pretrained embedding called \textbf{Word2Vec}, much like the continuous bag-of-words embedding of chapter 14. The objective turned \textbf{correlation} relationships between words into \textbf{distance} relationships in an embedding space.\par}\vskip 4pt
{\fontsize{9.4}{12.6}\selectfont\color{fg2}The resulting space did more than capture similarity. It showed a kind of emergent \textit{word arithmetic}:\par}\vskip 4pt
\begin{iquote}
\quotetext{V(king) $-$ V(man) + V(woman) $\approx$ V(queen)}\quotecite{A gender vector nobody trained for}
\end{iquote}
{\fontsize{9.4}{12.6}\selectfont\color{fg2}There were dozens of such vectors --- a plural vector, a vector from wild animals to their closest pet equivalent. \textbf{The model had not been trained for any of this explicitly.}\par}\vskip 4pt
\end{frame}

\framekicker{Section 15.5}
\begin{frame}[shrink=25]{How much a 2013 toy and a 2024 model have in common}
{\fontsize{9.4}{12.6}\selectfont\color{fg2}On the surface, nothing could seem further from Word2Vec than a large pretrained Transformer that generates fluent language on any topic. And yet:\par}\vskip 4pt
\begin{itable}{>{\raggedright\arraybackslash}p{0.3196\textwidth}>{\raggedright\arraybackslash}p{0.3102\textwidth}>{\raggedright\arraybackslash}p{0.3102\textwidth}}
\thead{} & \thead{Word2Vec} & \thead{Transformer} \\ \midrule
Goal & Embed tokens in a vector space & Embed tokens in a vector space \\
Learning principle & Tokens that co-occur end up close & Tokens that co-occur end up close \\
Distance function & Cosine distance & Cosine distance \\
Dimensionality & \textasciitilde{}1,000 per word & 1,000 to 10,000 per word \\
\end{itable}
{\fontsize{9.4}{12.6}\selectfont\color{fg2}You might object: a Transformer is trained to predict missing words, not to group tokens in a space. How does that objective relate to maximising dot-products between co-occurring tokens? \hilite{The answer is the attention mechanism.}\par}\vskip 4pt
\end{frame}

\framekicker{Section 15.5 $\cdot$ the geometric account}
\begin{frame}[shrink=25]{Attention is embedding-space refinement}
\leadpar{Attention learns a new embedding space by \textbf{linearly recombining} embeddings from a prior space, weighted toward tokens that are already close.}
{\fontsize{9.4}{12.6}\selectfont\color{fg2}Because the weighting is a dot product, tokens with high similarity contribute most to each other's new representation. Over training, this \textbf{pulls already-close vectors together} --- turning correlation relationships into proximity relationships.\par}\vskip 4pt
\begin{center}
\includegraphics[width=0.92\linewidth,height=0.60\textheight,keepaspectratio]{figs/ch15-why-effective.tex.pdf}
\end{center}
\vskip 3pt{\color{ink3}\fontsize{7.4}{9.4}\selectfont A Transformer learns a series of incrementally refined embedding spaces, each recombining elements of the previous one.\par}
\end{frame}

\framekicker{Section 15.5 $\cdot$ two properties}
\begin{frame}[shrink=25]{Continuous, and interpolative}
\begin{minipage}[t]{0.494\textwidth}
\begin{icardaccent}
\cardh{Semantically continuous}
Moving a little in the embedding space changes the human-facing meaning only a little. Word2Vec's space had this property too.
\end{icardaccent}
\end{minipage}%
\hfill
\begin{minipage}[t]{0.494\textwidth}
\begin{icardaccent}
\cardh{Semantically interpolative}
The midpoint between two points represents the \textbf{intermediate meaning} --- a direct consequence of each space being built by interpolating vectors from the previous one.
\end{icardaccent}
\end{minipage}%
{\fontsize{9.4}{12.6}\selectfont\color{fg2}This is not entirely unlike the brain. The key learning principle there is \textbf{Hebbian learning} --- \textit{neurons that fire together, wire together}. Correlations between firing events become proximity in the network.\par}\vskip 4pt
\begin{iband}[signal]
The Transformer, Word2Vec, and the brain all turn correlation relationships into vector-proximity relationships. All three are \hilite{maps of a space of information}.
\end{iband}
\end{frame}

\framekicker{Section 15.5 $\cdot$ the difference in degree}
\begin{frame}[shrink=25]{Vector programs, not vector functions}
{\fontsize{9.4}{12.6}\selectfont\color{fg2}Word2Vec is to today's language models as a logistic regression on MNIST pixels is to a state-of-the-art vision model. It was not even a deep network --- a single shallow layer. Today's Transformers stack dozens of attention and feedforward layers with parameter counts in the billions.\par}\vskip 4pt
{\fontsize{9.4}{12.6}\selectfont\color{fg2}With that representation power, and a more refined autoregressive objective, we are no longer confined to \textbf{linear} transformations like a gender vector.\par}\vskip 4pt
\begin{itable}{>{\raggedright\arraybackslash}p{0.4324\textwidth}>{\raggedright\arraybackslash}p{0.5076\textwidth}}
\thead{Word2Vec could store} & \thead{A large Transformer can store} \\ \midrule
\inlinecode{plural(cat) -> cats} & \inlinecode{write\_this\_in\_style\_of\_shakespeare("...poem...")} \\
\inlinecode{male\_to\_female(king) -> queen} & and \textbf{millions} of such programs \\
\end{itable}
{\fontsize{9.4}{12.6}\selectfont\color{fg2}These are so complex that it is more accurate to call them \textbf{vector programs} than vector functions --- highly nonlinear maps of the latent space onto itself. Not Python programs: no symbolic statements, no step-by-step data processing.\par}\vskip 4pt
\end{frame}

\framekicker{Section 15.5 $\cdot$ the database analogy}
\begin{frame}[shrink=25]{A database you can retrieve more from than you put in}
{\fontsize{9.4}{12.6}\selectfont\color{fg2}You can see a Transformer as a database: it stores information you retrieve via the tokens you pass in. Two differences matter enormously.\par}\vskip 4pt
\begin{minipage}[t]{0.494\textwidth}
\begin{icardaccent}
\cardh{It is continuous}
Not discrete entries but a \textbf{curve}. You can move along it to nearby points and interpolate between them --- so you can retrieve much more than was put in.
\end{icardaccent}
\end{minipage}%
\hfill
\begin{minipage}[t]{0.494\textwidth}
\begin{icardaccent}
\cardh{It stores programs}
RoBERTa holds facts, places, people, dates, relationships --- but \textbf{primarily} it is a database of vector programs.
\end{icardaccent}
\end{minipage}%
\begin{iband}[rose]
Interpolation is the same mechanism behind both outcomes: it leads to \textbf{generalization}, and it leads to \textbf{hallucination}. Not all of what you retrieve will be accurate or meaningful.
\end{iband}
\end{frame}
\note{This is the single most useful sentence in the chapter for a professional audience. Hallucination is not a bug to be patched out; it is the same property that makes the model useful, observed from the other side.}

\framekicker{Section 15.5 $\cdot$ the honest limit}
\begin{frame}[shrink=25]{Interpolation is not synthesis}
{\fontsize{9.4}{12.6}\selectfont\color{fg2}In the next chapter we push these models to billions of parameters and trillions of words. Their output can feel like magic --- like an intelligent operator sitting inside, pulling the strings.\par}\vskip 4pt
\begin{iquote}
\quotetext{These models are \textbf{fundamentally interpolative}. Thanks to attention, they learn an interpolative embedding space for a significant chunk of all text written in English. Wandering that space leads to interesting, unexpected generalizations --- but it cannot synthesize something fundamentally new with anything close to genuine, human-level intelligence.}\quotecite{Section 15.5, closing}
\end{iquote}
{\fontsize{9.4}{12.6}\selectfont\color{fg2}Keep this claim in mind through chapters 16 and 19, where it is tested against much larger models and much stronger claims.\par}\vskip 4pt
\end{frame}

\framekicker{Common failure modes}
\begin{frame}[shrink=25]{Four ways Transformer code goes quietly wrong}
\begin{minipage}[t]{0.494\textwidth}
\begin{icardbad}
\cardh{No causal mask on the decoder}
Training accuracy shoots up, generation produces nothing coherent. The model is reading its own label. \inlinecode{use\_causal\_mask=True}.
\end{icardbad}
\end{minipage}%
\hfill
\begin{minipage}[t]{0.494\textwidth}
\begin{icardbad}
\cardh{No positional embedding}
The model trains, converges, and scores several points too low with no error message. It is treating your sentences as \textbf{bags of words}.
\end{icardbad}
\end{minipage}%


\vskip 4pt

\begin{minipage}[t]{0.494\textwidth}
\begin{icardwarn}
\cardh{Tokenizer mismatched to weights}
Fine-tuning a pretrained backbone with the wrong tokenizer produces near-random performance. Load both from the \textbf{same preset}.
\end{icardwarn}
\end{minipage}%
\hfill
\begin{minipage}[t]{0.494\textwidth}
\begin{icardwarn}
\cardh{Default learning rate on fine-tuning}
Adam's default is \textasciitilde{}20$\times$ too large for a pretrained backbone and will wash out the representations in the first few batches.
\end{icardwarn}
\end{minipage}%
{\fontsize{9.4}{12.6}\selectfont\color{fg2}Three of these four produce \textbf{no exception and no warning} --- only a number that is worse than it should be. That is what makes them worth memorising.\par}\vskip 4pt
\end{frame}

\framekicker{Summary}
\begin{frame}[shrink=25]{What has to survive this chapter}
\begin{isteps}
\item \textbf{A language model learns p(token | past tokens)} --- and generates text of any length by being called in a loop, output becoming input.
\item \textbf{A masked language model} learns p(token | surrounding tokens), which is better for representing text than for generating it.
\item \textbf{Seq2seq} splits into an encoder that represents the source and a decoder that predicts the target one token at a time.
\item \textbf{Attention} scores by dot product, normalises by softmax, and returns a weighted sum --- giving direct access to any position in a sequence.
\end{isteps}
\end{frame}

\framekicker{Summary}
\begin{frame}[shrink=25]{What has to survive this chapter (2 of 2)}
\begin{isteps}
\item \textbf{The Transformer block} is attention (mixes positions, no nonlinearity) plus feedforward (mixes features, has the nonlinearity), with residuals and layer normalization around both.
\item \textbf{Attention is order-blind}, so positional embeddings are not optional --- leaving them out costs several points and raises no error.
\item \textbf{Pretraining then fine-tuning} beat training from scratch decisively: 93.7\% in one epoch against a 90\% ceiling.
\item \textbf{Attention turns correlation into proximity}, building interpolative embedding spaces --- the source of both generalization and hallucination.
\end{isteps}
\vskip 4pt
\reslink{NOTEBOOK}{03\_attention\_from\_scratch.ipynb}{../../course-slides/notebooks/ch15/03\_attention\_from\_scratch.ipynb}
\reslink{PAPER}{Vaswani et al., Attention Is All You Need}{https://arxiv.org/abs/1706.03762}
\reslink{NEXT}{Chapter 16 --- Text generation}{../ch16/index.html}
\vskip 2pt
\end{frame}

\end{document}
